% !TeX root = ../main.tex
% ---------------------------------------------------------------------------
% Stereographic Projection and Structural Analysis of Underground Excavations
%
% Promoted from an appendix to a main-narrative chapter: stereographic
% projection is a primary underground design tool, not slope-specific
% background. Roof and sidewall wedges, block removability, intersection
% geometry and excavation orientation all depend on it.
%
% Maps to: Hoek & Brown (1980) Ch. 4 "Graphical presentation of geological
% data" and Ch. 7 (structurally controlled instability);
% Brady & Brown (2006) Ch. 3 (projection tools) and Ch. 9
% "Excavation design in blocky rock".
% Net construction procedures are in Appendix~\ref{app:stereonet}.
% Related scripts: ../Stereonet and LEM.py, ../Kinematic Slope.py
% ---------------------------------------------------------------------------
\chapter{Stereographic Projection and Structural Analysis of Underground
Excavations}
\label{ch:structural-analysis}


\section{Projection in underground design}
\label{sec:projection-in-underground-design}

\subsection{Projection is an underground design tool}
\label{sec:projection-underground}
\scopenote{Stereographic projection is often introduced through slope
stability and then assumed to belong there. It does not. Structurally
controlled failure is one of the three primary underground failure modes
(Chapter~\ref{ch:empirical-support}), and every roof wedge, sidewall block,
portal failure and intersection collapse is identified by the same
construction. This chapter establishes the technique and applies it
underground; the slope-specific kinematic tests are not repeated here.}

\subsection{Where structural control governs underground}
\label{sec:structural-control-cases}
\scopenote{Crown and sidewall wedges in tunnels, roof blocks in large-span
caverns, wedges at junctions and cross-passages, portal and shaft-collar
blocks, and blocks released by a single fault in otherwise competent rock.
Cases where the rock mass is good by classification and still fails.}

\section{Projection fundamentals}
\label{sec:projection-fundamentals}

\subsection{The reference sphere and projection methods}
\label{sec:reference-sphere}
\scopenote{The lower hemisphere construction, equal-angle (Wulff) versus
equal-area (Schmidt) projection, and which to use for geometric construction
as against statistical contouring.}

\subsection{Plotting planes, poles and lines}
\label{sec:plotting-elements}
\scopenote{Great circles and poles for planes, points for lines. The
dip/dip-direction and strike/dip conventions, and the sign and hemisphere
errors that follow from mixing them. Step-by-step constructions in
Appendix~\ref{app:stereonet}.}

\subsection{Lines of intersection}
\label{sec:intersection-lines}
\scopenote{The intersection of two planes as a point on the net, its plunge
and trend, and why it is the single most important construction for
underground work: the intersection line is the direction a wedge slides.}

\section{Discontinuity data and its limitations}
\label{sec:discontinuity-data-and-its-limitations}

\subsection{Identifying discontinuity sets}
\label{sec:identifying-sets}
\scopenote{Pole density contouring, distinguishing genuine sets from
scatter and from random joints, mean set orientation, and dispersion measures
such as the Fisher constant.}

\subsection{Sampling bias and the Terzaghi correction}
\label{sec:terzaghi-correction}
\scopenote{Discontinuities subparallel to a borehole or scanline are
systematically under-sampled. The Terzaghi weighting correction, its blow-up
near the sampling direction, and the practical need for boreholes or scanlines
in more than one orientation. Connects to
Section~\ref{sec:orientation-bias}.}

\subsection{Errors in structural data processing}
\label{sec:projection-errors}
\scopenote{The error catalogue in Hoek and Brown: convention confusion,
uncorrected magnetic declination, merging incompatible data sets, contouring
too few poles, and treating a mean orientation as if it had no scatter.}

\section{From discontinuity sets to structural domains}
\label{sec:structural-domains}
\scopenote{Dividing an alignment into domains of consistent structure, and
why a single stereonet for a whole project hides the variation that governs
design.}

\section{Block and wedge formation}
\label{sec:block-and-wedge-formation}

\subsection{Block removability}
\label{sec:block-removability}
\scopenote{The geometric question that precedes any stability calculation:
given the discontinuity sets and the excavation surface, can the block move
into the opening at all? Finiteness and removability, and the joint pyramid
and excavation pyramid construction.}

\subsection{Roof, sidewall and face wedges}
\label{sec:tunnel-wedges}
\scopenote{Forming the tetrahedral wedges released by three intersecting sets
in the crown, each sidewall and the face of a tunnel; how the same joint sets
produce different critical wedges on differently oriented surfaces.}

\subsection{Wedge geometry, apex height and weight}
\label{sec:wedge-geometry}
\scopenote{Extracting apex height, base area and volume from the projection
and the excavation dimensions; the scaling of wedge weight with span that
makes large caverns structurally demanding.}

\section{Failure mode and excavation orientation}
\label{sec:failure-mode-and-excavation-orientation}

\subsection{Failure modes: free fall, single-plane and two-plane sliding}
\label{sec:wedge-modes}
\scopenote{Determining the mode from the position of the intersection line
relative to the excavation surface, and the different resisting mechanisms
each mode mobilises.}

\subsection{Excavation orientation as a design variable}
\label{sec:orientation-design}
\scopenote{Rotating the excavation on the stereonet to test how alignment
changes wedge size and mode. One of the few structural problems where a design
change removes the hazard rather than supporting it. Interacts with the
stress-driven orientation argument in Section~\ref{sec:cavern-orientation}.}

\subsection{Intersections, portals and shafts}
\label{sec:structural-intersections}
\scopenote{Additional free surfaces at junctions and cross-passages release
blocks that neither drive would form alone; the near-surface stress relief
that makes portal and shaft-collar blocks unusually mobile.}

\section{From geometry to support}
\label{sec:from-geometry-to-support}

\subsection{From geometry to stability and support}
\label{sec:geometry-to-support}
\scopenote{Passing wedge geometry to a limit equilibrium check, the role of
clamping stress and joint shear strength, and the required bolt force and
orientation. Design treatment in
Chapters~\ref{ch:ground-type-design} and~\ref{ch:support-interaction};
calculation sequence in Appendix~\ref{app:calculations}.}

\subsection{Probabilistic treatment of orientation}
\label{sec:probabilistic-orientation}
\scopenote{Discontinuity orientation is scattered, so wedge geometry is a
distribution rather than a single case. Sampling orientations to obtain a
distribution of wedge sizes and factors of safety.}

\section{Software and worked example}
\label{sec:software-and-worked-example}

\subsection{Software and digital workflows}
\label{sec:structural-software}
\scopenote{UnWedge and equivalent block analysis tools, photogrammetric and
LiDAR discontinuity extraction feeding projection directly, and the repository
scripts \code{Stereonet and LEM.py} and \code{Kinematic Slope.py}.}

\subsection{Worked example: cavern crown and sidewall wedges}
\label{sec:structural-worked-example}
\scopenote{Snowy 2.0 power station complex: joint sets from the geological
profile, wedge identification for the crown and each sidewall, and support
demand. Follows the UnWedge assessment referenced in
Chapter~\ref{ch:caverns}.}
